\documentclass[11pt,a4paper]{article}
\usepackage{amsmath,amssymb,amsthm,mathtools}
\usepackage[margin=1.1in]{geometry}
\usepackage{hyperref}
\usepackage{tikz}

\theoremstyle{plain}
\newtheorem{theorem}{Theorem}[section]
\newtheorem{proposition}[theorem]{Proposition}
\newtheorem{lemma}[theorem]{Lemma}
\newtheorem{corollary}[theorem]{Corollary}
\newtheorem{conjecture}[theorem]{Conjecture}
\theoremstyle{definition}
\newtheorem{definition}[theorem]{Definition}
\newtheorem{example}[theorem]{Example}
\newtheorem{remark}[theorem]{Remark}

\newcommand{\HH}{\mathrm{H}}
\newcommand{\N}{\mathbb{N}}
\newcommand{\tw}{\uparrow\uparrow}
\newcommand{\ev}{\mathrm{ev}}
\newcommand{\Tam}{\mathcal{T}}

%% ─────────────────────────────────────────────────────────────
%% DRAFT MARKERS.  To produce the submission version, change the
%% next line to \draftfalse ; nothing else needs to be touched.
\newif\ifdraft \drafttrue
\newcommand{\unseen}{\ifdraft\textsuperscript{\dag}\fi}
\newcommand{\draftnote}[1]{\ifdraft\par\medskip\noindent\fbox{\parbox{0.95\linewidth}{\small #1}}\par\medskip\fi}
%% ─────────────────────────────────────────────────────────────

\title{Bracketing hyperoperations:\\ a rank dichotomy for the Tamari order}
\author{Toshihisa Urahigashi}
\date{}

\begin{document}
\maketitle

\begin{abstract}
Let $\HH_r$ denote the Goodstein hyperoperation hierarchy, so that
$\HH_1$ is addition, $\HH_2$ multiplication, $\HH_3$ exponentiation and
$\HH_4$ tetration. Since $\HH_r$ is non-associative for $r\ge 3$, the value of a
bracketed expression $\HH_r(a,\dots,a)$ depends on the bracketing, i.e.\ on a
binary tree. We study how the value varies over the Tamari lattice of such trees.

We prove that a \emph{semi-associative law} runs through the whole hierarchy: for
every rank $r\ge3$ and all $x,y\ge2$,
\[
  \HH_r(\HH_r(x,y),z)\ \le\ \HH_r(x,\HH_r(y,z)),
\]
so a right rotation at any node never decreases the value and evaluation is
order-preserving from the Tamari lattice to $(\N,\le)$; the left comb is a
minimiser and the right comb, whose value is $\HH_{r+1}(a,n)$, a maximiser.

Our main result is that the \emph{degeneracy} of this law disappears at rank $4$.
We classify equality completely. For $x,y\ge2$:
\[
  r=3:\quad \text{equality}\iff z=1\ \text{ or }\ (y,z)=(2,2);
  \qquad
  r\ge4:\quad \text{equality}\iff z=1 .
\]
Thus rank $3$ carries exactly one non-trivial collision --- already visible for
$n=4$ leaves, where the Tamari cover $((aa)a)a\lessdot(aa)(aa)$ has equal
endpoints --- and from rank $4$ on the evaluation is strictly increasing along
every such cover. The mechanism is identified: the collapsing law
$(x^{y})^{z}=x^{yz}$, which produces the rank-$3$ equality, has no analogue at
rank~$4$ or above.

The result is not an asymptotic triviality. The two sides compare a growing
\emph{first} argument against a growing \emph{second} one, and small arguments do
not blow up with the rank: $\HH_r(2,2)=4$ for every $r\ge2$. Moreover the
hypothesis $y\ge2$ is sharp at every rank: for $y=1$ the inequality reverses.

All statements are formalized and machine-checked in Lean~4; the development
contains no \texttt{sorry} and every theorem is audited with \verb|#print axioms|.
\end{abstract}

\section{Introduction}

Exponentiation is not associative, and it is classical that among all bracketings
of $a^{a^{\cdots^{a}}}$ with $n$ occurrences of $a\ge2$ the right-associated one
is the largest; its value is the tetration $a\tw n$. What appears not to have been
examined is the \emph{order} structure behind this fact, and its behaviour further up the
hyperoperation hierarchy.

Two strands of existing work bracket the question. On the one hand, the
\emph{associative spectrum} of a binary operation, introduced by Cs\'ak\'any and
Waldhauser~\cite{CsW}, measures non-associativity by counting how many distinct
values the $C_{n-1}$ bracketings can take; they prove in particular that
exponentiation is \emph{Catalan}, i.e.\ distinct bracketings induce distinct
operations. That line of work counts values. On the other hand, the set of
bracketings carries the Tamari lattice~\cite{Tamari}, generated by right
rotations, whose proof-theoretic content was isolated by Zeilberger~\cite{Zeil}
as a semi-associative law. That line of work orders bracketings. The present
paper connects the two: we show that for hyperoperations the Tamari order is
\emph{respected} by evaluation, and we locate precisely where the correspondence
degenerates.

\subsection*{Results}

Throughout, $\HH_r$ is the hyperoperation hierarchy fixed in
Section~\ref{sec:prelim}, with $\HH_1(a,b)=a+b$, $\HH_2(a,b)=ab$,
$\HH_3(a,b)=a^b$ and $\HH_4(a,b)=a\tw b$.

\begin{theorem}[Rank $3$, complete solution]\label{thm:A}
Let $x\ge1$ and $y\ge2$. Then $(x^{y})^{z}\le x^{(y^{z})}$ for every $z\ge0$.
If moreover $x\ge2$, equality holds if and only if $z=1$ or $(y,z)=(2,2)$.
The hypothesis $y\ge2$ cannot be dropped: for $y=1$ the inequality fails.
\end{theorem}

\begin{theorem}[Rank $4$]\label{thm:B}
Let $x\ge2$ and $y\ge2$. Then $(x\tw y)\tw z\le x\tw(y\tw z)$ for every $z\ge0$.
If moreover $z\ge2$, the inequality is strict.
\end{theorem}

\begin{proposition}[Tamari monotonicity]\label{prop:tamari}
Let $a\ge2$ and $r\ge3$. If $T\lessdot U$ is a Tamari cover, i.e.\ $U$ arises from
$T$ by a right rotation at some node, then $\ev_{r,a}(T)\le\ev_{r,a}(U)$.
Consequently $\ev_{r,a}$ is order-preserving for the Tamari order. Moreover, for the
right comb $R_n$ with $n$ leaves,
$\ev_{r,a}(R_n)=\HH_{r+1}(a,n)$.
\end{proposition}

\begin{corollary}[Dichotomy]\label{cor:C}
The monotonicity of Proposition~\ref{prop:tamari} is \emph{not strict} at rank~$3$:
the cover $((aa)a)a\lessdot(aa)(aa)$ carries equal values, both $256$, when $a=2$.
For every $r\ge4$, by contrast, a root rotation is strict as soon as the value of the
third subtree is at least $2$, which holds automatically for $a\ge2$.
\end{corollary}

\begin{theorem}[All ranks]\label{thm:D}
For every $r\ge3$ and all $x\ge2$, $y\ge2$ and $z\ge0$,
\[
  \HH_r\bigl(\HH_r(x,y),z\bigr)\ \le\ \HH_r\bigl(x,\HH_r(y,z)\bigr).
\]
\end{theorem}

\begin{theorem}[Complete equality classification]\label{thm:E}
Let $x\ge2$ and $y\ge2$.
\begin{enumerate}
\item For every $r\ge2$, equality holds in Theorem~\ref{thm:D} when $z=1$.
\item For $r=3$ the only further equality is $(y,z)=(2,2)$.
\item For $r\ge4$ and $z\ge2$ the inequality is \emph{strict}.
\end{enumerate}
Hence for $r\ge4$ equality holds if and only if $z=1$.
\end{theorem}

Theorems~\ref{thm:D} and~\ref{thm:E} are proved in Section~\ref{sec:general} by
inductions on the rank whose engine is the identity-like inequality \eqref{eq:SUM}
below; the case $r=3$ is Theorem~\ref{thm:A} and the case $r=4$ can also be
obtained directly, as in Section~\ref{sec:rank4}.

\begin{remark}[The statement is not an asymptotic triviality]\label{rem:nontriv}
Two features are worth isolating. First, \eqref{eq:HT} does not follow from
monotonicity alone: on the left the \emph{first} argument grows,
$x\mapsto\HH_r(x,y)>x$, while on the right the \emph{second} argument grows,
$z\mapsto\HH_r(y,z)$, so the two coordinates compete. Second, high rank does not
make everything large: since $\HH_r(a,1)=a$ for $r\ge2$, one gets
\[
  \HH_r(2,2)=\HH_{r-1}(2,\HH_r(2,1))=\HH_{r-1}(2,2)=\dots=\HH_2(2,2)=4
\]
for \emph{every} $r\ge2$. A degenerate corner therefore survives at every rank, and
it is exactly this corner that carries the rank-$3$ equality of
Theorem~\ref{thm:E}(2) and that has to be treated separately in the proof of
Theorem~\ref{thm:E}(3).
\end{remark}

\subsection*{Formalization}

Every theorem stated here is proved in Lean~4 on top of Mathlib. The development
defines the hyperoperation by the usual double recursion and \emph{verifies the
convention} by proving $\HH_1=+$, $\HH_2=\cdot$, $\HH_3=(\ )^{(\ )}$; this matters
because the literature on hyperoperations uses several incompatible indexings and
initial values. Section~\ref{sec:formal} describes the development.

\section{Preliminaries}\label{sec:prelim}

\begin{definition}[Hyperoperation~\cite{Good}]\label{def:hyper}
Define $\HH:\N\times\N\times\N\to\N$ by
\[
\HH_0(a,b)=b+1,\quad
\HH_1(a,0)=a,\quad
\HH_2(a,0)=0,\quad
\HH_{r+3}(a,0)=1,
\]
\[
\HH_{r+1}(a,b+1)=\HH_{r}\bigl(a,\HH_{r+1}(a,b)\bigr).
\]
\end{definition}

\begin{lemma}\label{lem:conv}
For all $a,b$ we have $\HH_1(a,b)=a+b$, $\HH_2(a,b)=ab$ and $\HH_3(a,b)=a^{b}$.
Moreover $\HH_r(a,1)=a$ for $r\ge2$.
\end{lemma}

We write $a\tw b:=\HH_4(a,b)$, so $a\tw 0=1$ and $a\tw(b+1)=a^{(a\tw b)}$.

\begin{definition}[Bracketings and evaluation]
Let $\Tam_n$ denote the set of binary trees with $n$ leaves. For $r,a\in\N$ define
$\ev_{r,a}:\Tam_n\to\N$ by $\ev_{r,a}(\text{leaf})=a$ and
$\ev_{r,a}(T\cdot U)=\HH_r\bigl(\ev_{r,a}(T),\ev_{r,a}(U)\bigr)$.
\end{definition}

\begin{definition}[Right rotation]
The \emph{right rotation at the root} sends $((A\cdot B)\cdot C)$ to
$(A\cdot(B\cdot C))$. The Tamari order on $\Tam_n$ is generated by right
rotations performed at arbitrary nodes; that this order is a lattice, with the left
comb as least and the right comb as greatest element, is the theorem of Friedman and
Tamari~\cite{FT}, whose title already names the relation a \emph{loi
demi-associative}~\cite{Tamari,Zeil}.
\end{definition}

Thus the inequality
\begin{equation}\label{eq:HT}
  \HH_r\bigl(\HH_r(x,y),z\bigr)\ \le\ \HH_r\bigl(x,\HH_r(y,z)\bigr)
  \tag{HT$_r$}
\end{equation}
is exactly the assertion that a right rotation, applied with $x,y,z$ the values of
the three subtrees, does not decrease the value.

\section{Rank three}\label{sec:rank3}

At rank $3$ the collapsing law $(x^{y})^{z}=x^{yz}$ reduces \eqref{eq:HT} to a
comparison of exponents.

\begin{lemma}\label{lem:mul-le-pow}
If $y\ge2$ then $yz\le y^{z}$ for every $z$. If moreover $z\ge3$ the inequality is strict.
\end{lemma}

\begin{proposition}\label{prop:mul-eq-pow}
If $y\ge2$ then $yz=y^{z}$ if and only if $z=1$ or $(y,z)=(2,2)$.
\end{proposition}

\begin{proof}
If $z=0$ then $yz=0\ne1=y^{z}$. If $z=1$ both sides are $y$. If $z=2$ then
$yz=y^{z}$ reads $2y=y^{2}$, i.e.\ $y=2$ after cancelling $y>0$. If $z\ge3$ the
inequality is strict by Lemma~\ref{lem:mul-le-pow}. The converse is immediate.
\end{proof}

\begin{proof}[Proof of Theorem~\ref{thm:A}]
Since $(x^{y})^{z}=x^{yz}$, the inequality follows from
Lemma~\ref{lem:mul-le-pow} and monotonicity of $t\mapsto x^{t}$ for $x\ge1$.
For $x\ge2$ the map $t\mapsto x^{t}$ is injective, so equality holds iff
$yz=y^{z}$, and Proposition~\ref{prop:mul-eq-pow} applies. For the last claim,
$(x^{1})^{z}=x^{z}$ while $x^{(1^{z})}=x$, and $x^{z}>x$ whenever $x,z\ge2$.
\end{proof}

\begin{example}[Failure of strictness]\label{ex:collide}
Take $a=2$ and $n=4$. The five bracketings evaluate under $\HH_3$ to
\[
  ((aa)a)a=256,\quad (a(aa))a=256,\quad (aa)(aa)=256,\quad
  a((aa)a)=65536,\quad a(a(aa))=65536 .
\]
In particular the Tamari cover $((aa)a)a\lessdot(aa)(aa)$ has equal endpoints,
so $\ev_{3,2}$ is monotone but not strictly monotone; it is also not injective.
\end{example}

\section{Rank four}\label{sec:rank4}

At rank $4$ no collapsing law is available and a different argument is required.
The proof below is by induction on $z$; the base case $z=1$ must be treated
separately, since the inductive hypothesis at $z=0$ is too weak to survive the step.

\begin{lemma}[Squaring]\label{lem:sq}
If $x\ge2$ and $m\ge1$ then $(x\tw m)^{2}\le x\tw(m+1)$.
\end{lemma}

\begin{proof}
Induction on $m$. For $m=1$ the claim is $x^{2}\le x^{x}$, which holds as $x\ge2$.
For the step put $t=x\tw m$; then
$(x\tw(m+1))^{2}=(x^{t})^{2}=x^{2t}\le x^{(x^{t})}=x\tw(m+2)$,
where $2t\le 2^{t}\le x^{t}$ uses $t\ge1$ and $x\ge2$.
\end{proof}

\begin{remark}
The direct inequality $t^{2}\le x^{t}$ is \emph{false} in general --- take $x=2$,
$t=3$ --- so the inductive formulation in Lemma~\ref{lem:sq} is essential.
\end{remark}

\begin{lemma}\label{lem:add-two}
If $y\ge2$ and $m\ge2$ then $m+2\le y^{m}$. If moreover $m\ge3$ then $m+3\le y^{m}$.
\end{lemma}

\begin{proof}[Proof of Theorem~\ref{thm:B}]
Write $y=y'+1$ with $y'\ge1$ and induct on $z$.
For $z=0$ the left side is $1$ and the right side is $x\ge1$.
For $z=1$ both sides equal $x\tw y$.
Assume $z\ge1$ and set $m:=y\tw z$, so $m\ge y\ge2$ and $y'\le m$.
By Lemma~\ref{lem:add-two} we have $y^{m}\ge m+2$, so $y^{m}=k+1$ with $k\ge m+1$.
Writing $A:=x\tw y=x^{(x\tw y')}$ we compute
\begin{align*}
  (x\tw y)\tw(z+1)
  &= A^{\,(x\tw y)\tw z}
   \ \le\ A^{\,x\tw m}
   \ =\ \bigl(x^{x\tw y'}\bigr)^{x\tw m}
   \ =\ x^{(x\tw y')(x\tw m)}\\
  &\le x^{(x\tw m)^{2}}
   \ \le\ x^{\,x\tw(m+1)}
   \ \le\ x^{\,x\tw k}
   \ =\ x\tw(k+1)
   \ =\ x\tw\bigl(y\tw(z+1)\bigr),
\end{align*}
using the inductive hypothesis, $x\tw y'\le x\tw m$, Lemma~\ref{lem:sq}, and
$m+1\le k$. If $z\ge2$ then $m\ge4$, so $y^{m}\ge m+3$ by
Lemma~\ref{lem:add-two} and hence $m+1<k$; the penultimate step is then strict.
\end{proof}

\begin{remark}\label{rem:strict}
The argument above gives strictness whenever $y\tw z\ge3$, which covers all $z\ge3$
and all $z\ge2$ with $y\ge3$. The remaining case is $y=z=2$ with $x\ge2$ arbitrary,
and there strictness holds too:
\[
  (x\tw2)\tw2=(x^{x})^{x^{x}}=x^{\,x^{x+1}},
  \qquad
  x\tw(2\tw2)=x\tw4=x^{\,x^{x^{x}}},
\]
and $x+1<x^{x}$ for $x\ge2$. Hence the inequality of Theorem~\ref{thm:B} is strict
for every $z\ge2$; only the uniform argument degrades at $y\tw z=2$.
\end{remark}

\section{The dichotomy}\label{sec:dicho}

\begin{proof}[Proof of Corollary~\ref{cor:C}]
Monotonicity along a root rotation is exactly \eqref{eq:HT} instantiated at the
values of the three subtrees; since $\ev_{r,a}(T)\ge2$ for $a\ge2$, the
hypotheses of Theorems~\ref{thm:A} and~\ref{thm:B} are met.
Non-strictness at rank~$3$ is Example~\ref{ex:collide}. Strictness at rank~$4$
is Theorem~\ref{thm:B}.
\end{proof}

\begin{proof}[Proof of Proposition~\ref{prop:tamari}]
Induction on the derivation of the cover. At the root the claim is
Theorem~\ref{thm:D} instantiated at the values of the three subtrees, which are all
$\ge2$ because $a\ge2$. If the rotation takes place inside the left subtree, one uses
monotonicity of $\HH_r(\cdot,v)$ in its first argument; inside the right subtree,
monotonicity of $\HH_r(u,\cdot)$ in its second argument, valid for $u\ge2$. Both
monotonicity statements hold for every $r\ge1$ and are proved by a double induction
on the rank and on the argument. The formula for the right comb is immediate from
$\HH_{r+1}(a,n+1)=\HH_r(a,\HH_{r+1}(a,n))$ together with $\HH_{r+1}(a,1)=a$.
\end{proof}

The mechanism is worth stating plainly. At rank $3$ the identity
$(x^{y})^{z}=x^{yz}$ turns \eqref{eq:HT} into $yz\le y^{z}$, an inequality which
is \emph{tight}: it becomes an equality at $(y,z)=(2,2)$, and this single
coincidence is the source of every collapse in Example~\ref{ex:collide}. At rank
$4$ there is no such identity, the comparison is genuinely two-dimensional, and
the slack is enormous.

\section{General rank}\label{sec:general}

We record the general statement and the shape of a proof.

We now prove Theorem~\ref{thm:D}. The engine is the following inequality; write
\begin{equation}\label{eq:SUM}
  \HH_{r-1}\bigl(\HH_r(x,a),\HH_r(x,b)\bigr)\ \le\ \HH_r(x,\,a+b).
  \tag{SUM$_r$}
\end{equation}

\begin{proposition}\label{prop:reduction}
Let $r\ge3$ and $x\ge2$.
\begin{enumerate}
\item If \textup{(HT$_{r-1}$)} holds, then \textup{(SUM$_r$)} holds for all $a,b\ge1$.
\item If \textup{(SUM$_r$)} holds and $u+v\le\HH_{r-1}(u,v)$ for all $u,v\ge2$,
then \textup{(HT$_r$)} holds for all $y\ge2$.
\end{enumerate}
\end{proposition}

\begin{proof}
(1) Induct on $a$. For $a=1$ both sides are equal:
\[
\HH_{r-1}\bigl(\HH_r(x,1),\HH_r(x,b)\bigr)
  =\HH_{r-1}\bigl(x,\HH_r(x,b)\bigr)=\HH_r(x,b+1).
\]
For the step, $\HH_r(x,a+1)=\HH_{r-1}(x,\HH_r(x,a))$, so by (HT$_{r-1}$)
\[
\HH_{r-1}\bigl(\HH_{r-1}(x,\HH_r(x,a)),\HH_r(x,b)\bigr)
\le \HH_{r-1}\bigl(x,\HH_{r-1}(\HH_r(x,a),\HH_r(x,b))\bigr),
\]
and the inductive hypothesis bounds the inner term by $\HH_r(x,a+b)$, giving
$\HH_{r-1}(x,\HH_r(x,a+b))=\HH_r(x,a+b+1)$.

(2) Induct on $z$, with $z=0,1$ as in Theorem~\ref{thm:B}. For the step set
$m:=\HH_r(y,z)\ge2$; then
$\HH_r(\HH_r(x,y),z+1)=\HH_{r-1}(\HH_r(x,y),\HH_r(\HH_r(x,y),z))
\le\HH_{r-1}(\HH_r(x,y),\HH_r(x,m))\le\HH_r(x,y+m)\le\HH_r(x,\HH_{r-1}(y,m))
=\HH_r(x,\HH_r(y,z+1))$.
\end{proof}

\begin{proof}[Proof of Theorem~\ref{thm:D}]
Induction on $r$. The base case $r=3$ is Theorem~\ref{thm:A}. For the step, the
inductive hypothesis is exactly (HT$_{r-1}$), so Proposition~\ref{prop:reduction}(1)
gives (SUM$_r$), and Proposition~\ref{prop:reduction}(2) then gives (HT$_r$), the
hypothesis $u+v\le\HH_{r-1}(u,v)$ for $u,v\ge2$ being Lemma~\ref{lem:add}.
\end{proof}

\begin{lemma}\label{lem:add}
For $u,v\ge2$ and $s\ge2$ we have $u+v\le\HH_s(u,v)$.
\end{lemma}

\begin{proof}
$u+v\le uv=\HH_2(u,v)$, and $\HH_2(u,v)\le\HH_s(u,v)$ by monotonicity in the rank.
\end{proof}

\subsection*{Strictness}

\begin{lemma}\label{lem:addstrict}
For $u,v\ge2$ with $u\ge3$ or $v\ge3$, and $s\ge2$, we have $u+v<\HH_s(u,v)$.
Equality $u+v=\HH_s(u,v)$ occurs exactly at $u=v=2$, where both sides are $4$.
\end{lemma}

\begin{proof}
$u+v<uv$ holds iff $(u-1)(v-1)>1$, i.e.\ iff $u\ge3$ or $v\ge3$; and
$uv=\HH_2(u,v)\le\HH_s(u,v)$. At $u=v=2$ we have $u+v=4=\HH_2(2,2)=\HH_s(2,2)$ by
Remark~\ref{rem:nontriv}.
\end{proof}

\begin{proof}[Proof of Theorem~\ref{thm:E}]
(1) $\HH_r(\HH_r(x,y),1)=\HH_r(x,y)=\HH_r(x,\HH_r(y,1))$ since $\HH_r(a,1)=a$.
(2) is Theorem~\ref{thm:A}.
For (3) we induct on $r\ge4$, the base $r=4$ being Theorem~\ref{thm:B} with
Remark~\ref{rem:strict}. Let $r\ge5$ and write $z=w+1$ with $w\ge1$, and
$m:=\HH_r(y,w)\ge2$.

If $y\ge3$ or $w\ge2$, then $y\ge3$ or $m\ge4$, so Lemma~\ref{lem:addstrict} makes
the last step of the proof of Proposition~\ref{prop:reduction}(2) strict, and
strictness propagates because $\HH_r(x,\cdot)$ is strictly increasing.

The remaining case is $y=z=2$. Put $A:=\HH_r(x,2)=\HH_{r-1}(x,x)\ge4$. Then
\[
  \HH_r(\HH_r(x,2),2)=\HH_{r-1}(A,A),
  \qquad
  \HH_r(x,\HH_r(2,2))=\HH_r(x,4)=\HH_{r-1}\bigl(x,\HH_{r-1}(x,A)\bigr),
\]
using $\HH_r(2,2)=4$ from Remark~\ref{rem:nontriv}. Since $A=\HH_{r-1}(x,x)$, the
required inequality
\[
\HH_{r-1}\bigl(\HH_{r-1}(x,x),A\bigr)<\HH_{r-1}\bigl(x,\HH_{r-1}(x,A)\bigr)
\]
is the inductive hypothesis at rank $r-1$ applied to $(x,x,A)$, and $A\ge2$.
\end{proof}

The auxiliary facts used above --- monotonicity of $\HH_r(x,\cdot)$ in its second
argument, monotonicity in the rank, and $n+1\le\HH_r(x,n)$ for $x\ge2$, $r\ge2$,
$n\ge1$ --- are elementary and are proved in the formalization by a sequence of
inductions.

The hypotheses of Theorem~\ref{thm:D} are sharp. For $y=1$ the inequality fails at
every rank $r\ge3$: then $\HH_r(x,1)=x$, so the left side is $\HH_r(x,z)$, whereas
$\HH_r(1,z)=1$ makes the right side $\HH_r(x,1)=x$, and $\HH_r(x,z)>x$ as soon as
$z\ge2$. For $x=0$ the operations degenerate.

\section{The Lean development}\label{sec:formal}

The development is written for Lean~4 (v4.30.0) with Mathlib (v4.30.0) and is
available at
\begin{center}
\url{https://github.com/turahigashi/hyper-tamari}
\end{center}
It consists of three files, \verb|HyperTamari/Basic.lean|,
\verb|HyperTamari/Tetration.lean| and \verb|HyperTamari/General.lean|, together with
the raw \verb|lake build| output recording the axiom audit
(\verb|logs/axiom-audit.txt|). It contains, among others:
\begin{itemize}
\item \verb|hyper|, the operation of Definition~\ref{def:hyper}, together with
      \verb|hyper_one|, \verb|hyper_two|, \verb|hyper_three| certifying
      Lemma~\ref{lem:conv};
\item \verb|mul_eq_pow_iff| and \verb|ht_rank3_eq_iff|, Theorem~\ref{thm:A};
\item \verb|BTree|, \verb|eval| and the cover relation \verb|Rot| (right rotation at
      an arbitrary node), with \verb|Rot.eval_le_rank3| and \verb|Rot.eval_le_rank4|,
      Proposition~\ref{prop:tamari};
\item \verb|eval_rightComb|, the closed form for the right comb;
\item \verb|rank3_rotation_not_strict|, Example~\ref{ex:collide};
\item \verb|mul_le_pow_of_two_le| and \verb|mul_lt_pow_of_three_le|,
      Lemma~\ref{lem:mul-le-pow};
\item \verb|sq_tet_le|, Lemma~\ref{lem:sq}, and \verb|add_two_le_pow|,
      \verb|add_three_le_pow|, Lemma~\ref{lem:add-two};
\item \verb|tet_ht| and \verb|tet_ht_strict_of_three_le|, Theorem~\ref{thm:B};
\item \verb|rank3_vs_rank4_dichotomy|, Corollary~\ref{cor:C};
\item \verb|sum_lemma| and \verb|ht_general|, Proposition~\ref{prop:reduction} and
      Theorem~\ref{thm:D}, together with the auxiliary
      \verb|succ_le_hyper|, \verb|hyper_mono_arg|, \verb|hyper_rank_mono| and
      \verb|add_le_hyper|;
\item \verb|hyper_mono_base| and \verb|Rot.eval_le|, the general-rank form of
      Proposition~\ref{prop:tamari};
\item \verb|tet_ht_strict_of_two_le|, the strengthening in Remark~\ref{rem:strict};
\item \verb|hyper_two_two| ($\HH_r(2,2)=4$) and \verb|hyper_one_base|,
      Remark~\ref{rem:nontriv};
\item \verb|add_lt_hyper|, \verb|ht_strict_general| and
      \verb|equality_classification|, Theorem~\ref{thm:E}.
\end{itemize}
Every numbered statement of Sections~\ref{sec:prelim}--\ref{sec:general} has a named
counterpart in the development; the list above gives the correspondence in full.

\draftnote{\textbf{Draft note --- not part of the submission.}
The artifact has \emph{not} yet been deposited in a long-term archive, so no DOI is
quoted above. A version DOI will be minted and inserted here at the moment this
paper is submitted, and the archive record will be linked to the submission.
Until then the repository URL is the only pointer.}
No declaration uses \verb|sorry|, \verb|native_decide| or any private axiom.
Every one of the \textbf{77} theorems in these three files is audited with
\verb|#print axioms|, and the audit is machine-checked to cover them all
(declared 77, audited 77, missing 0); \verb|sorryAx| does not occur.
The axiom dependencies are \verb|propext| alone (2 theorems),
\verb|propext, Quot.sound| (45) and
\verb|propext, Classical.choice, Quot.sound| (30); the uses of choice come from
library lemmas on \verb|Nat| and are not essential to the arguments.

\section{Related work}\label{sec:related}

Cs\'ak\'any and Waldhauser~\cite{CsW} introduce the associative spectrum and show
that exponentiation is Catalan: for pairwise distinct prime leaves, distinct
bracketings give distinct values. Their concern is the number of values; the
order in which those values arrange themselves is not addressed, and the Tamari
lattice appears in their bibliography only as the source of the enumeration of
bracketings. Huang and Lehtonen~\cite{HL} extend the invariant to the
\emph{associative-commutative spectrum} and determine it for exponentiation
(their Proposition 6.1.5): the associative spectrum of exponentiation is $C_{n-1}$
and its ac-spectrum is $n^{n-1}$. Their question is again how many distinct term
operations the bracketings produce; ours is how the resulting values are
\emph{arranged}, and neither the Tamari order nor any monotonicity statement occurs
in that line of work. Huang, Mickey and Xu~\cite{HMX} compute the associative
spectrum of a further operation. Tamari~\cite{Tamari} enumerated the bracketings, and Friedman and
Tamari~\cite{FT} proved that the right-rotation order on them is a finite lattice ---
their title already calls the relation a \emph{loi demi-associative}, which is
precisely the law we show to run through the whole hyperoperation hierarchy.
Zeilberger~\cite{Zeil} gives a sequent calculus for that law. Gryszka~\cite{Gry} studies the comparison of power towers
by normalizing them, which concerns right-associated towers rather than general
bracketings.

\section*{Use of AI assistance}

In accordance with current disclosure practice, the author states the following.

Artificial-intelligence assistants were used substantially in the preparation of
this work. The question studied here was formulated by the author. Large language
models were then used to search for the shape of the proofs, to write the entire
Lean~4 development, to draft this text, and to check the literature; further
language models were used as adversarial readers, and their criticism led to the
present form of the main theorem, in particular to the classification of the
equality cases in Theorem~\ref{thm:E}. All mathematical decisions, the choice of
what to claim, and the final text are the author's responsibility.

The reader is asked to note what this does and does not mean. Every mathematical
assertion in Sections~\ref{sec:rank3}--\ref{sec:general} has been formalised and
checked by the Lean~4 kernel, and the accompanying development is distributed with
the paper; the correctness of those statements therefore does not rest on the
reliability of any language model. What does rest on human and machine judgement,
and is stated here so that it can be checked independently, is the choice of
definitions, the claim of novelty, and the attribution of prior work.

\draftnote{\textbf{Draft note --- not part of the submission.}
Items marked \textsuperscript{\dag} in the list below have \emph{not} been consulted
in the original at the time of writing; the bibliographic data were taken from
secondary sources and independently cross-checked against Crossref, but the
contents were not read. They are:
\textbf{[FT]} Friedman--Tamari 1967, cited only for the origin of the lattice and of
the term \emph{loi demi-associative} (the phrase occurs in the title itself);
\textbf{[Gry]} Gryszka 2022, cited only as related work on comparing power towers;
\textbf{[Tamari]} Tamari 1962, cited only for the enumeration of bracketings and for
the two combs being the extreme elements.
None of the results of the present paper depends on these three items: they are
cited for attribution and context, and every mathematical statement used in the
proofs is either proved here or formalised in the accompanying Lean development.
Before submission each of the three must be checked in the original, and this note
removed by setting \texttt{\textbackslash draftfalse}.}

\begin{thebibliography}{9}
\bibitem{CsW} B. Cs\'ak\'any and T. Waldhauser,
  \emph{Associative spectra of binary operations},
  Multiple-Valued Logic \textbf{5} (2000), 175--200; arXiv:1102.2124.
\bibitem{FT}\unseen{}  H. Friedman and D. Tamari,
  \emph{Probl\`emes d'associativit\'e: une structure de treillis finis induite par une
  loi demi-associative},
  Journal of Combinatorial Theory \textbf{2} (1967), no.~3, 215--242.
\bibitem{Good} R. L. Goodstein,
  \emph{Transfinite ordinals in recursive number theory},
  Journal of Symbolic Logic \textbf{12} (1947), 123--129.
\bibitem{Gry}\unseen{}  K. Gryszka, \emph{How to compare power towers?},
  Lithuanian Mathematical Journal \textbf{62} (2022), no.~2, 192--206.
\bibitem{HL} J. Huang and E. Lehtonen,
  \emph{The associative-commutative spectrum of a binary operation},
  Discrete Mathematics \textbf{346} (2023), 113468; arXiv:2202.11826.
\bibitem{HMX} J. Huang, M. Mickey and J. Xu,
  \emph{The nonassociativity of the double minus operation},
  arXiv:1710.05651 (2017).
\bibitem{Tamari}\unseen{}  D. Tamari,
  \emph{The algebra of bracketings and their enumeration},
  Nieuw Archief voor Wiskunde (3) \textbf{10} (1962), 131--146.
\bibitem{Zeil} N. Zeilberger,
  \emph{A sequent calculus for a semi-associative law},
  Logical Methods in Computer Science \textbf{15} (2019), no.~1, 9:1--9:23.
\end{thebibliography}

\end{document}
